% What Minimal Einstein--Cartan Torsion Does for the Bounce and Cannot Do
% for Dark Energy
% Paper 1N -- gr-qc / Classical and Quantum Gravity Note
% Merges paper1a_ech_nogo.tex (v1A.0.127) and paper1c_nogo_survey/main.tex
% (v1C.0.16) into a single publishable Note, per
% project-context/PORTFOLIO_DECISION_2026-09-02.md Sec.3 Track B and Addendum.
% Houston Golden -- 2026
% arXiv submission: gr-qc (cross-list astro-ph.CO, hep-th)
%
% Provenance: this Note supersedes paper1a_ech_nogo.tex and
% paper1c_nogo_survey/main.tex as the single publication target for the
% closed ECH dark-energy line. The two source papers remain archived
% (P1A: Zenodo doi:10.5281/zenodo.21481838; P1C: repository history) and
% are not further submitted independently.
% Compile: pdflatex main && bibtex main && pdflatex main && pdflatex main

\documentclass[aps,prd,twocolumn,superscriptaddress,preprintnumbers,nofootinbib,longbibliography,floatfix]{revtex4-2}

\usepackage{amsmath,amssymb,amsfonts}
\usepackage{graphicx}
\usepackage{bm}
\usepackage{hyperref}
\usepackage{xcolor}
\usepackage{booktabs}
\usepackage{multirow}
\usepackage{dcolumn}
\usepackage{enumitem}
\usepackage{mathtools}
\usepackage[utf8]{inputenc}
\usepackage{float}

\hypersetup{
  colorlinks=true,
  linkcolor=blue!60!black,
  citecolor=blue!60!black,
  urlcolor=blue!60!black,
  pdfnewwindow=true
}

\newcommand{\Leff}{\Lambda_{\rm eff}}
\newcommand{\Lobs}{\Lambda_{\rm obs}}
\newcommand{\rhocrit}{\rho_{\rm crit}}
\newcommand{\rhoPl}{\rho_{\rm Pl}}
\newcommand{\MPl}{M_{\rm Pl}}
\newcommand{\ie}{i.e.}
\newcommand{\eg}{e.g.}
\newcommand{\cf}{cf.}
\newcommand{\etal}{\textit{et~al.}}
\newcommand{\fnl}{f_{\rm NL}}
\newcommand{\artifact}[1]{\href{https://github.com/Hubify-Projects/bigbounce/blob/main/#1}{repository artifact}}
\newcommand{\paperVersion}{v1N.0.1}
\newcommand{\paperTimestamp}{September 2, 2026}

\begin{document}

\title{What Minimal Einstein--Cartan Torsion Does for the Bounce and Cannot
Do for Dark Energy}

\author{Houston Golden}
\email{houston@hubify.com}
\affiliation{Independent Researcher, Los Angeles, California, USA}

\date{\paperTimestamp\ (\paperVersion)}

\begin{abstract}
Pop{\l}awski's Einstein--Cartan black-hole cosmology replaces the classical
singularity with a torsion-supported bounce: eliminating the
non-propagating spin connection from the minimal Einstein--Cartan--Holst
(ECH) action generates an effective spin--spin four-fermion contact term
that is repulsive at high spin density, the same mechanism that produces
Pop{\l}awski's torsion-induced avoidance of gravitational
singularities~\cite{Poplawski2010,Poplawski2011}. We ask whether the
identical algebraic mechanism can also source the observed late-time
dark-energy density and answer, systematically, no. Consolidating two
companion papers into a single Note, we (i)~derive the minimal
axial--axial contact interaction
$-(3\kappa/16)[\gamma^2/(1+\gamma^2)]J_5^2$ from the same Cartan equation
that produces the bounce term, and show its declared direct-channel,
hard-cutoff, standard mean-field NJL scalar projection is repulsive,
$G_s=-3\kappa/16$, so the real homogeneous gap equation for this
condensate channel has no nonzero solution; (ii)~prove, for canonical
scalar matter, a perturbation-transparency theorem: torsion vanishes at
every classical perturbation order and the Holst sector decouples
identically from the scalar and tensor equations of motion, because the
Holst dual contraction vanishes pointwise by the algebraic Bianchi
identity on the torsion-free branch; and (iii)~catalog fourteen distinct
mechanism-class constraints -- amplitude suppressions, naturalness
deficits, decoupling identities, and classification arguments spanning
seven structural foundations and six observational branches -- that
jointly close the four enumerated channels (NJL contact, one-loop Holst
correction, Immirzi running, parity-odd CMB coupling) by which minimal
ECH could plausibly connect its bounce-scale torsion to a late-time
$\Lambda$-like density. A six-member spanning list of dimension-four
parity-odd local densities admitted by the minimal field content is
shown to have rank four as a set of operators and rank two modulo total
derivatives -- a generating set, not an independent basis -- with every
member either an exact total derivative or a Fierz-closed,
$\MPl^{-2}$-suppressed contact term; the on-shell ECH torsion itself
carries both an axial and a trace-vector irrep at finite Barbero--Immirzi
parameter $\gamma$ (ratio $\beta/\alpha=1/2\gamma$; the tensor irrep
vanishes identically), a correction to the purely-axial reading used in
an earlier catalog draft. The result is a clean structural dichotomy: the
same contact term that plausibly supports Pop{\l}awski's bounce mechanism
is, under minimal coupling, parity-even, Planck-suppressed, and
classically transparent to perturbations -- exactly the properties that
make it unable to generate the late-time acceleration independently
attributed to dark energy. No ECH dark-energy or birefringence prediction
is made; this is a channel-level, not operator-level, closure within the
stated minimal-coupling scope.
\end{abstract}

\keywords{Einstein--Cartan gravity, Holst action, torsion bounce,
Pop{\l}awski cosmology, four-fermion interaction, dark energy, no-go}

\maketitle

%% ============================================================
\section{Introduction}\label{sec:intro}

Minimal Einstein--Cartan--Holst (ECH) gravity promotes the spin connection
to an independent field sourced algebraically by fermion spin currents. In
loop-quantized cosmology the resulting torsion term replaces the classical
singularity with a quantum bounce, and in Pop{\l}awski's
black-hole-cosmology programme the same torsion term is the mechanism by
which a collapsing fermionic interior avoids a singularity and instead
bounces, producing a new expanding universe inside what an outside
observer sees as a black hole~\cite{Poplawski2010,Poplawski2012,Poplawski2016,UngerPoplawski2019}.
The physical content of that mechanism is a spin--spin four-fermion
contact interaction: eliminating the non-propagating Cartan connection in
favor of the fermionic spin current generates an effective interaction
that grows with spin density and becomes repulsive well before the
classical curvature singularity is reached, halting collapse. This is the
same contact term this Note derives below, and the same sign.

Because the ECH bounce and Pop{\l}awski's proposed link from bounce-scale
torsion to a late-time vacuum energy~\cite{Poplawski2011,Poplawski2012}
are built from the same minimal field content, a natural question is
whether the contact term responsible for the bounce can \emph{also} source
the dark-energy density observed today, $\rho_{\Lambda,\rm
obs}\approx(2.25\,\text{meV})^4$. This Note answers that question
systematically for the minimal-coupling ECH action: it merges the
algebraic derivation and rigorous transparency theorem of an earlier
companion paper~\cite{Golden2026P1a} with the systematic no-go survey of a
second companion~\cite{Golden2026P1cArxiv} into a single, self-contained
Classical and Quantum Gravity Note. Sec.~\ref{sec:theory} derives the
contact term from the same Cartan equation that produces Pop{\l}awski's
bounce interaction and states its conditional-sign gap-equation result.
Sec.~\ref{sec:transparency} proves the perturbation-transparency theorem,
the Note's sole Tier-I rigorous result. Sec.~\ref{sec:barriers} catalogs
the fourteen mechanism-class constraints that jointly close the four
enumerated dark-energy channels; Sec.~\ref{sec:routes} summarizes the two
amplitude-budget route closures; Sec.~\ref{sec:completeness} states the
corrected operator-list argument. Sec.~\ref{sec:discussion} draws the
positive/negative dichotomy this Note is organized around, and
Sec.~\ref{sec:conclusions} concludes.

This is a channel-level assessment of the minimal-coupling theory, not an
operator-level completeness theorem, and it makes no claim about
non-minimal completions, propagating torsion, a dynamical Immirzi field,
or quantum loops beyond the stated scope of each result.

%% ============================================================
\section{The Contact Term: Same Mechanism as the Bounce}\label{sec:theory}

We use the first-order ECH action with constant Barbero--Immirzi parameter
$\gamma$ and a minimally coupled Dirac field,
\begin{align}\label{eq:ECH}
S[e,\omega,\psi] ={}& \frac{1}{4\kappa}\int
\left(\epsilon_{IJKL}+\frac{2}{\gamma}\eta_{I[K}\eta_{L]J}\right)\nonumber\\
&\hspace{1.0cm}\times e^I\wedge e^J\wedge R^{KL}(\omega)
+ S_{\rm D}[e,\omega,\psi],
\end{align}
$\kappa=8\pi G=8\pi/\MPl^2$, mostly-plus signature
$\eta_{IJ}=\mathrm{diag}(-,+,+,+)$, $\epsilon_{0123}=+1$. There is no
independent torsion-squared term in Eq.~\eqref{eq:ECH}; any four-fermion
expression below is an on-shell result obtained only after the connection
equation is solved. Varying $\omega_{IJ}$ gives the algebraic equation
$\mathcal Q_\gamma(e^{[I}\wedge T^{J]})=\mathcal J^{IJ}$ with
$\mathcal Q_\gamma=\star+\gamma^{-1}\mathbb 1$ invertible for every real
finite nonzero $\gamma$~\cite{Freidel2005}. For minimal fermion coupling
the spin current is dual to $J_5^I=\bar\psi\gamma^I\gamma^5\psi$, and
solving the connection equation and back-substituting
(Freidel--Minic--Takeuchi~\cite{Freidel2005}, Eqs.~17 and 23, in the
normalization bridge $4\pi G=\kappa/2$) gives the minimal axial--axial
contact interaction
\begin{equation}\label{eq:4fermi}
\mathcal L_{4\psi}=-\frac{3\kappa}{16}\frac{\gamma^2}{1+\gamma^2}\,
(J_5^I J_{5I}).
\end{equation}
This is Pop{\l}awski's torsion-bounce interaction restricted to minimal
coupling: it is generated by the identical elimination of the
non-propagating connection, grows quadratically with the fermionic spin
density, and carries the same repulsive sign that halts gravitational
collapse in the black-hole-cosmology mechanism~\cite{Poplawski2010}. The
Einstein--Cartan coefficient ($\gamma\to\infty$) is the maximal magnitude,
since $\gamma^2/(1+\gamma^2)<1$.

At finite $\gamma$ the connection solution
$e_I{}^\mu C_{\mu JK}\propto\left(\tfrac12\epsilon_{IJKL}J_5^L
-\gamma^{-1}\eta_{I[J}J^5_{K]}\right)$ carries two irreducible pieces: an
axial part with coefficient $\alpha\propto\gamma^2/(\gamma^2+1)$ and a
trace-vector part with coefficient $\beta\propto\gamma/(\gamma^2+1)$,
ratio $\beta/\alpha=1/(2\gamma)$; the tensor irrep vanishes identically
for minimal coupling. This corrects an earlier catalog draft's
``purely-axial'' reading of the on-shell torsion, which holds only in the
strict Einstein--Cartan limit $\gamma\to\infty$, not at the finite
physical $\gamma$ used elsewhere in this programme
(theory-audit record~\cite{TorsionOnshell2026}).

\subsection{Finite-density benchmark and the NJL sign result}

A deliberately elevated homogeneous normalization illustrates only the
scale of Eq.~\eqref{eq:4fermi}:
\begin{equation}
\frac{\kappa n_\psi^2}{\rho_\Lambda}\simeq
3.6\times10^{-69}\left(\frac{n_\psi}{100\,\mathrm{cm}^{-3}}\right)^2,
\end{equation}
${\approx}68$ orders of magnitude below $\rho_{\Lambda,\rm obs}$ even at
an ISM-like number density. This coefficient-one benchmark omits the
actual $3/16$ and finite-$\gamma$ factors; number density also does not
fix the renormalized composite $\langle J_5^I J_{5I}\rangle$, a vacuum
stress tensor, or an equation of state.

In the declared direct-channel, hard-four-momentum-cutoff, standard
mean-field NJL convention, Fierz rearrangement of Eq.~\eqref{eq:4fermi}
projects onto the scalar channel $G_s(\bar\psi\psi)^2$ with
\begin{equation}
G_s=-\frac{3\kappa}{16}<0,
\end{equation}
repulsive. The real homogeneous scalar gap equation of this convention
therefore has no nonzero solution: this specific truncation admits no
condensate. This conditional-sign result does not exclude other
truncations, species structures, non-minimal couplings, or propagating
torsion; the full regulated derivation is carried in the archived
companion~\cite{Golden2026P1a}.

%% ============================================================
\section{Perturbation Transparency}\label{sec:transparency}

The catalog's sole Tier-I rigorous result is the following. Consider the
classical ECH action with an invertible tetrad, canonical scalar matter,
and real finite nonzero constant $\gamma$, on a local patch with matched
background, initial, and boundary data (standard falloff, vanishing
first-order boundary variation). Solve the algebraic connection equation
before expanding in perturbations. Then the reduced action on the
zero-spin branch is the Einstein--scalar action, so the scalar equations
and tensor evolution operators coincide with GR at every perturbative
order for matched data. Equality of right- and left-helicity
\emph{solutions} additionally requires parity-symmetric initial data. The
complex singular values $\gamma=\pm i$, nontrivial global/topological
sectors, quantum loops or anomalies, fermion sources, non-minimal matter,
a dynamical Immirzi field, and propagating torsion are outside this
statement.

\emph{Proof.} (1)~A canonical scalar has zero spin density. (2)~The zero
source and invertible $\mathcal Q_\gamma$ give $e^{[I}\wedge T^{J]}=0$;
contracting twice with the dual frame shows the kernel is trivial for an
invertible tetrad, so $T^I=0$. (3)~The connection reduces to
Levi-Civita, $\Gamma^\lambda{}_{\mu\nu}=\mathring\Gamma^\lambda{}_{\mu\nu}$.
(4)~The Holst term
$\tfrac12\epsilon^{\mu\nu\rho\sigma}R_{\mu\nu\rho\sigma}(\mathring\Gamma)$
vanishes identically on a torsion-free connection by the first (algebraic)
Bianchi identity $R_{\mu[\nu\rho\sigma]}=0$, which holds for any
torsionless connection independent of metric compatibility. This
Bianchi-vanishing is a single-curvature identity, distinct from the
two-curvature Pontryagin density $R\widetilde R$; no total-derivative
argument is load-bearing. On an FRW background the tensor evolution
operators are identical for both helicities,
$\mathcal E_R=\mathcal E_L=\partial_\eta^2+2\mathcal H\partial_\eta+k^2$,
so minimal ECH generates no GW birefringence, tensor chirality, or
$TB$/$EB$ CMB cross-power from parity-symmetric initial data, and the
cubic action for $\zeta$ receives zero contribution from the Holst term
in the stated domain. The theorem is an on-connection-shell equality of
local classical reduced actions -- not an off-shell equality of the
original first-order ECH and Einstein--scalar actions -- and its
all-orders content follows from the algebraic Cartan constraint and the
pointwise Bianchi identity, not an order-by-order calculation. A full
self-contained statement and second-order Holst-term verification is
carried in the archived companion~\cite{Golden2026P1a}.

%% ============================================================
\section{The Barrier Catalog}\label{sec:barriers}

Seven structural ``foundations'' (parameter-space, symmetry, and
naturalness arguments; Barriers 1--7) and six observational branches
(CMB, gravitational-wave, and relic-abundance channels; Barriers 8--14,
with Branch~H carrying two logically independent entries B8 and B14)
together give fourteen distinct mechanism-class constraints that jointly
close the four enumerated dark-energy channels (Table~\ref{tab:barriers}).
`Distinct' and `mechanism-class' mean no barrier is a logical consequence
of another and each probes a separate physical failure mode; this does
not assert disjoint assumptions, nor that each is an independently
rigorous no-go theorem. Several barriers (B5, B6, B7, B10, B13) are
general naturalness or classification arguments rather than sharp
ECH-specific calculations, and B9 is an explicitly heuristic closure
conditional on stated equilibrium assumptions that realistic quantum
bounces can violate; B9 and B14 are never used as stand-alone closures of
any route. The sole Tier-I leg among the fourteen is B14 (perturbation
transparency, Sec.~\ref{sec:transparency}), and it is Tier-I strictly on
the zero-spin branch its hypotheses cover.

\begin{table}[tb]
\caption{\label{tab:barriers} The fourteen catalog entries -- fourteen
distinct mechanism-class constraints -- on minimal-ECH dark-energy
routes. B8 and B14 close the same observable channel (tensor parity
violation) on disjoint matter branches and are logically independent
(Sec.~\ref{sec:transparency}).}
\centering
\footnotesize
\begin{ruledtabular}
\begin{tabular}{clll}
\# & Barrier & Src.\ & Blocks \\
\midrule
1 & Mass-Coupling Lock & A & Torsion as DE \\
2 & Topological-Shift Duality & B & Pseudoscalar mass prot. \\
3 & Scalar-Tensor Universality & C & Distinctive FRW geom. \\
4 & Planck Suppression & D & Disformal coupling \\
5 & Scale Separation & E & Global vacuum coupling \\
6 & Attractor Sensitivity & F & IC transfer to DE \\
7 & Parameter Immunity & G & Cyclic vacuum select. \\
8 & Parity-Even Interaction & H & Tensor chirality (ferm.) \\
9 & Liouville Conservation & J & Reversible selection \\
10 & UV$\to$IR Specificity & L & Generic/specific bridge \\
11 & Decoupling Universality & L/M & Light gauge coupling \\
12 & Vacuum Amplif.\ Ceiling & M & GW amplitude \\
13 & Gravitational Democracy & N/O & Relics, baryogenesis \\
14 & Perturbation Transparency & H & Tensor chirality (scal.) \\
\end{tabular}
\end{ruledtabular}
\normalsize
\end{table}

\medskip\noindent\textbf{B1--B2 [R1, Founds.\ A--B].} Ultralight torsion
modes require $g_{\rm eff}\sim H_0/\MPl\sim10^{-61}$; reaching $g_{\rm
eff}\sim1$ needs $m_T\sim\MPl$, the standard CC hierarchy
$\sim10^{-122}$. In metric-affine gravity, mass protection $\iff$ no
geometric fingerprint: configurations protecting the pseudoscalar mass
eliminate the geometric content and vice versa.

\medskip\noindent\textbf{B3--B4 [R2, Founds.\ C--D].} Diffeomorphism
invariance forces the torsion fluctuation to couple like any other FRW
scalar, with no ECH-specific observable. Disformal torsion couplings are
Planck-suppressed, $\mathcal O(10^{-122})$ at $H_0$.

\medskip\noindent\textbf{B5--B7 [R3, Founds.\ E--G]; general naturalness
arguments.} No mechanism transfers the global vacuum integral across
$N_{\rm tot}\approx92$--94 $e$-folds; an inflationary attractor washes out
bounce initial data while sensitivity to them destabilizes inflation;
$\gamma$ is fixed by the LQG area spectrum with no landscape for cyclic
selection.

\medskip\noindent\textbf{B8 [R1, Branch H].} $(J_5)^2$ is parity-\emph{even}
(product of two axial currents), so it cannot generate tensor chirality.
Logically independent of B14: B14's hypotheses require zero spin density,
precisely the condition B8's nonzero $J^5$ violates.

\medskip\noindent\textbf{B9 [R2, Branch J]; heuristic.} Phase-space
conservation prevents irreversible post-bounce state selection under
closed non-dissipative evolution with no particle production; scenarios
that break these assumptions evade B9 and are constrained instead by
B10--B11, so B9 is never used as a stand-alone closure.

\medskip\noindent\textbf{B10--B11 [R3, Branch L/M]; B10 general
naturalness.} A Planck-to-$H_0$ bridge must be either generic or
bounce-specific, never both; gauge fields decouple from the
Planck-suppressed torsion sector universally, so ECH cannot preferentially
couple to photons or dark-energy degrees of freedom without new
non-minimal couplings.

\medskip\noindent\textbf{B12 [R3, Branch M].} Bounce-epoch
$\Omega_{\rm GW}^{\rm ECH}\lesssim(\rhocrit/\rhoPl)^2\simeq0.07$--$0.17$
using the LQG-bounce window; an order-of-magnitude ceiling, not directly
comparable to the nHz PTA spectral density without propagating through
the transfer function (deferred).

\medskip\noindent\textbf{B13 [R4, Branch N/O]; structural observation.}
Torsion couples democratically to all spin-$\tfrac12$ species and cannot
preferentially source baryogenesis or relic abundances without
species-dependent non-minimal couplings; logically independent of B11 (gauge
universality) and B4 (coupling magnitude).

\medskip\noindent\textbf{B14 [R2--R4, zero-spin branch, Branch H].} Sole
Tier-I leg; Sec.~\ref{sec:transparency}. Because its hypotheses assume
zero spin density, it does not bear on R1; what it establishes for
R2--R4 is that their classical zero-spin baseline is exactly inert, so
any signal must come entirely from the quantum or non-minimal ingredient
defining each route -- bounded by Sec.~\ref{sec:routes}.

%% ============================================================
\section{Route 2--4: Amplitude and Naturalness Closures}\label{sec:routes}

Route~1 (NJL contact) is closed by Sec.~\ref{sec:theory}: amplitude
suppressed by $\MPl^{-2}$ and parity-even. Route~2 (one-loop graviton
corrections to the Holst term) is anchored in the published one-loop
renormalization of the Holst-plus-fermion sector~\cite{ShapiroTeixeira2014}
and closes against the observed birefringence amplitude with
${\gtrsim}58$ orders of suppression margin; its dark-energy leg is closed
operator-level for constant Nieh--Yan coefficient (Sec.~\ref{sec:completeness})
and only at the naturalness level for a dynamical one. Route~3 (Immirzi
running) is bounded by the integrated Benedetti--Speziale one-loop
flow~\cite{BenedettiSpeziale2011run}, leaving the contribution
$61$--$67$ orders below $\rho_{\Lambda,\rm obs}$. Route~4 (parity-odd CMB
coupling via a spectator ALP or the fermion axial current) is not
amplitude-suppressed: a free-normalization ALP fit reproduces both the
WMAP+Planck birefringence signal $\beta_{\rm obs}=0.342^\circ\pm0.094^\circ$~\cite{Eskilt2022,Minami2020}
(comparable to ACT DR6, $\beta=0.215^\circ\pm0.074^\circ$~\cite{DiegoPalazuelos2025})
\emph{and} $\rho_\Lambda$, but minimal ECH derives neither the required
ultralight mass $m_\theta\sim H_0$ nor the fitted coupling from first
principles -- it relocates rather than solves the cosmological-constant
problem, so R4 closes as a naturalness/explanatory-deficit objection, not
an amplitude exclusion.

\begin{table*}[tb]
\caption{\label{tab:evidentiary_status} Evidentiary status of each leg: (I)
rigorous within stated scope, (II) structural argument, (III)
ansatz-level dimensional estimate. No leg is claimed more strongly than
recorded here.}
\centering
\begin{ruledtabular}
\begin{tabular}{ll}
Leg & Status \\
\colrule
\parbox[t]{0.20\textwidth}{Perturbation\\transparency} &
\parbox[t]{0.65\textwidth}{\raggedright \textbf{(I)} Rigorous within
stated scope.} \\[4pt]
\parbox[t]{0.20\textwidth}{R1 (NJL contact)} &
\parbox[t]{0.65\textwidth}{\raggedright \textbf{(II)+(III)}: parity-even
structure is algebraic; $\MPl^{-2}$ suppression is an ansatz-level
estimate.} \\[4pt]
\parbox[t]{0.20\textwidth}{R2 (one-loop)} &
\parbox[t]{0.65\textwidth}{\raggedright \textbf{(III)} birefringence leg;
\textbf{(II)} dark-energy leg (both branches).} \\[4pt]
\parbox[t]{0.20\textwidth}{R3 (Immirzi running)} &
\parbox[t]{0.65\textwidth}{\raggedright \textbf{(II)+(III)}:
mass-dimension lock is structural; the chiral-count bound is a loose
Tier-III ceiling.} \\[4pt]
\parbox[t]{0.20\textwidth}{R4 (CMB/spectator-ALP)} &
\parbox[t]{0.65\textwidth}{\raggedright \textbf{(II)} structural
naturalness objection, not an amplitude exclusion.} \\
\end{tabular}
\end{ruledtabular}
\end{table*}

The R1--R3 legs achieve mathematical amplitude closure (suppressed by
explicit $\MPl$ powers and loop factors, many orders below observable
thresholds regardless of normalization); R4 achieves
naturalness/explanatory closure. These are logically distinct closure
modes, classified separately above. The phenomenological parameter
$\alpha/M$ that fits $\beta_{\rm obs}$ is therefore best understood as an
R4-bounded coupling with the dark-energy density supplied by an unrelated
sector, rather than by ECH-internal physics.

%% ============================================================
\section{The Operator List: Rank, Not Basis}\label{sec:completeness}

The barrier catalog and route closures above bound the amplitude of
representative parity-odd operators. A referee may reasonably ask whether
the conclusion survives the full local operator space. Under the
construction rule of zero additional derivative order, parity-odd
$\varepsilon$ contraction, full index contraction to a scalar density,
and mass dimension exactly four, the admitted densities built from the
tetrad, the torsionful curvature, the algebraically-fixed torsion, and
the minimal axial current $J^5$ are spanned by a six-member list
$\{$O1--O6$\}$. Independent symbolic adjudication~\cite{OperatorBasisAdj2026}
establishes that this list has \textbf{rank four} as a set of densities
(two exact relations, $\mathrm{O1}=\mathrm{O6}$ and
$\mathrm{O1}=\tfrac12\mathrm{O4}-\mathrm{O2}$) and \textbf{rank two}
modulo total derivatives -- a deliberately redundant generating set of
recognizable invariants, \emph{not} a linearly independent basis; the
list is asserted to span the rule-admitted space from the construction
rule itself, not by exhaustive enumeration.

Every member reduces to one of three classes: an exact total derivative
(O2, the Nieh--Yan density; O3, the Pontryagin density; and O1=O6 on the
torsion-free branch), a Fierz-closed four-fermion contact term suppressed
by $\MPl^{-2}$ (O5, the surviving vacuum-energy content,
$\mathrm{O5}^{[4]}=-\tfrac32\kappa(J^5\!\cdot\!J^5)$), or identically
vanishing. On the on-shell ECH branch the single-curvature pair (O1, O6)
vanishes by the algebraic Bianchi identity on the torsion-free connection,
exactly as in Sec.~\ref{sec:transparency}; the $\varepsilon$-contracted
torsion-square O4 was reported vanishing under a purely-axial reading of
the on-shell torsion in an earlier catalog draft, but independent on-shell
symbolic evaluation~\cite{TorsionOnshell2026} shows the finite-$\gamma$
ECH torsion carries a genuine trace-vector irrep alongside the axial one
(Sec.~\ref{sec:theory}), so O4 is nonzero on the axial$\times$trace-vector
cross term and joins O5 in the Fierz-closed disposal class rather than
vanishing outright. This does not change the closure: O4's surviving
content is still $\MPl^{-2}$-suppressed and inside the same Fierz-closed
basis as O5, so Route~2's dark-energy leg for constant Nieh--Yan
coefficient is unaffected, only its stated mechanism is corrected.

%% ============================================================
\section{Discussion}\label{sec:discussion}

The two halves of this Note's title are the same calculation read in two
directions. Eliminating the non-propagating ECH connection in favor of
the fermionic spin current generates exactly one interaction --
Eq.~\eqref{eq:4fermi} -- and that interaction is simultaneously
(a)~repulsive at high spin density, which is the physical content of
Pop{\l}awski's torsion-avoided singularity and bounce mechanism, and
(b)~parity-even, Planck-suppressed, and (on the zero-spin branch)
classically transparent to every scalar and tensor perturbation --
exactly the properties a viable dark-energy source cannot have. Minimal
ECH is not a two-purpose mechanism with one purpose realized and the
other merely unproven; the barrier catalog of Sec.~\ref{sec:barriers}
shows that every distinct route from the bounce-scale torsion sector to a
late-time $\Lambda$-like density is closed by a separate, named failure
mode, and Sec.~\ref{sec:completeness} shows those failure modes are not
an artifact of an incomplete operator enumeration.

This closure is specific to the ECH mechanism class under minimal
coupling; it says nothing about non-minimal torsion couplings,
propagating torsion, a dynamical Immirzi field, or the bounce's
predicted galaxy-spin observable, which is tested independently by the
DESI-scale chirality survey~\cite{Golden2026P4}. Nor does it bear on
bounce cosmology generally: the matter-bounce $\fnl=-35/16$ signature and
NANOGrav-band gravitational-wave forecasts are properties of the
matter-bounce class, not of ECH torsion, and are pursued in a separate
research line~\cite{Golden2026P2}. What this Note closes is narrower and
more useful to state cleanly: the specific claim that the same minimal
spin-torsion sector producing the ECH bounce also supplies the observed
dark-energy density.

%% ============================================================
\section{Conclusions}\label{sec:conclusions}

Minimal Einstein--Cartan--Holst torsion, under strictly minimal fermion
coupling, generates one physical interaction: a spin--spin four-fermion
contact term, repulsive at high density, algebraically identical to the
mechanism underlying Pop{\l}awski's torsion-avoided singularity and
black-hole-cosmology bounce. This Note shows, at channel-amplitude
granularity and within a stated minimal-coupling scope, that the same
interaction cannot additionally source the observed late-time dark-energy
density: its NJL scalar projection admits no condensate in the declared
truncation (Sec.~\ref{sec:theory}), it decouples identically from
classical perturbations on the zero-spin branch by a rigorous
Bianchi-identity theorem (Sec.~\ref{sec:transparency}), fourteen
mechanism-class constraints jointly close all four enumerated dark-energy
channels (Sec.~\ref{sec:barriers}--\ref{sec:routes}), and a rank-four
spanning list of the admissible dimension-four parity-odd operators shows
this conclusion is not an artifact of an incomplete operator basis
(Sec.~\ref{sec:completeness}). The positive result -- what minimal ECH
does for the bounce -- and the negative result -- what it cannot do for
dark energy -- are two readings of the same algebraic elimination, and
this Note's contribution is to state both precisely, with the closure
scope and evidentiary strength of every claim recorded explicitly rather
than implied.

\subsection*{Data and Code Availability}
The symbolic verification scripts supporting the transparency theorem,
the on-shell torsion-irrep computation, and the operator-list adjudication
are publicly available in the companion papers' repository,
\url{https://github.com/Hubify-Projects/bigbounce}:
\begin{center}\footnotesize
\artifact{arxiv/scripts/dim4_parityodd_enumeration.py} \\
\artifact{research/theory_audit/operator_basis_adjudication_2026_08_07.py} \\
\artifact{research/theory_audit/ech_torsion_onshell_2026_08_08.py}
\end{center}
The earlier version of the algebraic torsion-elimination derivation and
its regulated NJL check is archived, with its own verification scripts,
as an immutable deposit at
\href{https://doi.org/10.5281/zenodo.21481838}{doi:10.5281/zenodo.21481838}
(CC-BY-4.0)~\cite{Golden2026P1a}. This Note supersedes that paper and the
companion no-go survey~\cite{Golden2026P1cArxiv} as the single submission
target for the closed ECH dark-energy line; both source manuscripts
remain archived and are not independently resubmitted.

\begin{acknowledgments}
This Note builds on the published spin-torsion cosmology of Nikodem
Pop\l{}awski~\cite{Poplawski2010,Poplawski2011,Poplawski2012,Poplawski2016},
on the published derivations of Simone Mercuri~\cite{Mercuri2009} and of
Laurent Freidel, Djordje Minic, and Tatsu Takeuchi~\cite{Freidel2005}
connecting the Barbero--Immirzi parameter to parity-violating fermion
interactions, and on the published one-loop renormalization results of
Ilya Shapiro and Poliane Teixeira~\cite{ShapiroTeixeira2014} and of Dario
Benedetti and Simone Speziale~\cite{BenedettiSpeziale2011run} that anchor
the Route-2/Route-3 budgets. None of the named researchers was involved
in this work, and no endorsement is implied.

\emph{AI-assisted methods.}---This work was carried out with an agentic
AI research pipeline operated under the author's direction, used for
systematic barrier-cataloging, symbolic and dimensional cross-checking,
literature triage, and manuscript preparation. All scientific claims,
derivations, and numerical results were verified by the author against
the committed computational artifacts in the public reproducibility
tree, which together with the repository git history constitutes a
public audit trail. The author takes sole responsibility for all
scientific claims, derivations, numerical results, and bibliographic
attributions in this paper. No external funding was received for this
research.
\end{acknowledgments}

\bibliographystyle{apsrev4-2}
\bibliography{references}

\end{document}
